\documentclass[11pt]{article}

\usepackage[letterpaper,margin=1in]{geometry}
\usepackage{amsmath,amssymb,amsthm}
\usepackage{booktabs,tabularx}
\usepackage{microtype}
\usepackage{xcolor}
\usepackage{xurl}
\usepackage[hidelinks]{hyperref}

\newtheorem{theorem}{Theorem}
\newtheorem{lemma}[theorem]{Lemma}
\newcommand{\C}{\mathbb C}
\newcommand{\Q}{\mathbb Q}
\newcommand{\cP}{\mathcal P}
\newcommand{\ii}{\mathrm i}
\newcommand{\Hess}{\operatorname{Hess}}

\title{An explicit Hessian-nilpotent quartic in 54 variables\\
witnessing the failure of Zhao's Vanishing Conjecture}
\author{Alec Kriebel\\[2pt]\large with heavy assistance from ChatGPT 5.6 Sol}
\date{First public release: 21 July 2026, 13:11:39 UTC\\
Superseded-status correction: 21 July 2026, 15:13:18 UTC}

\begin{document}
\maketitle

\begin{quote}
\textbf{Superseded status.} Exploration 03 replaces the 27-variable cubic and
54-variable quartic dimension counts by 22 and 44. Thompson had also posted a
24-variable cubic map before this note's first release. This document is
retained as a timestamped derivation of the factor-reusing 13-variable stable
model and its original 54-variable certificate; its compactness claim is
historical, not current.
\end{quote}

\begin{abstract}
We execute a factor-reusing Bass--Connell--Wright degree reduction of the
recently announced three-dimensional noninjective Keller map, obtaining a
cubic homogeneous noninjective Keller map in 27 variables.  Applying the de
Bondt--van den Essen symmetric reduction then gives an explicit homogeneous
quartic Hessian-nilpotent polynomial in 54 variables, with 598 monomials, whose
gradient Keller map is noninjective.  This is an explicit witness to the
failure of Zhao's Vanishing Conjecture, with the expanded polynomial and a
two-point collision supplied as machine-readable exact certificates.  The
underlying three-dimensional map was posted by Levent Alpöge on 20 July 2026,
crediting Akhil Mathew for posing the question and Claude Fable for producing
the example.
\end{abstract}

\noindent\textbf{Verification status.}
The algebra has passed the supplied exact symbolic checks, but this short note
has not been peer reviewed. Independent expert scrutiny is welcome.

\section{Statement of the result}

Let $\ii^2=-1$.  The straight-line construction below defines a cubic
homogeneous map $h:\C^{27}\to\C^{27}$.  For $A,B\in\C^{27}$, set
\begin{equation}
  \boxed{\cP(A,B)=\ii\sum_{j=1}^{27}h_j(A+\ii B)B_j.}
  \label{eq:potential}
\end{equation}

\begin{theorem}\label{thm:main}
The polynomial $\cP$ is homogeneous of degree four, has 598 monomials when
expanded, and its Hessian is nilpotent.  The polynomial map
\begin{equation}
  \Gamma:\C^{54}\longrightarrow\C^{54},\qquad
  \Gamma(Z)=Z-\nabla\cP(Z),
  \label{eq:gamma}
\end{equation}
has Jacobian determinant one and is noninjective.  Moreover,
\begin{equation}
  \Delta^m\cP^m=0\quad\text{for every }m\geq1,
  \qquad
  \Delta^m\cP^{m+1}\neq0\quad\text{for infinitely many }m\geq0.
  \label{eq:vc-counterexample}
\end{equation}
Thus $\cP$ is a counterexample to Zhao's Vanishing Conjecture.
\end{theorem}

The formulas below and ordinary differentiation form a finite exact
certificate.  The accompanying \texttt{potential\_sparse.json} is the fully
expanded polynomial over $\Q(\ii)$, and \texttt{collision.json} contains two
distinct points with the same image under~\eqref{eq:gamma}.

\section{A 13-variable cubic stable model}

Write $u=1+xy$.  Postcomposing the map announced in~\cite{alpoge} by the
inverse of its linear part gives the identity-linear Keller map
$\Phi=(\Phi_1,\Phi_2,\Phi_3)$:
\begin{align}
\Phi_1&=x-\frac32x^2y-\frac12x^3z,\nonumber\\
\Phi_2&=y+3xu^2z+3xy^2(4+3xy),\label{eq:phi}\\
\Phi_3&=u^3z+y^2u(4+3xy).\nonumber
\end{align}
Thus $\det J\Phi=1$, and the three points
\begin{equation}
(0,0,-\tfrac14),\qquad(1,-\tfrac32,\tfrac{13}2),\qquad
(-1,\tfrac32,\tfrac{13}2)
\label{eq:three-points}
\end{equation}
have the same image $(0,0,-\tfrac14)$.

We use the following stable degree-reduction gadget.  Suppose a term $cPQ$ is
to be removed from the $k$th coordinate of a current map $M$.  With new
variables $a,b$, replace
\begin{equation}
 M_k\longmapsto M_k-c(a+P)(b+Q),\qquad
 (a,b)\longmapsto(a+P,b+Q).
 \label{eq:two-gadget}
\end{equation}
This is a pre- and post-composition of $M\times\mathrm{id}_{\mathbb A^2}$ by
triangular automorphisms.  A previously created output factor may be reused:
with one new variable $a$, one may replace
\begin{equation}
 M_k\longmapsto M_k-c(a+P)M_\ell,\qquad a\longmapsto a+P.
 \label{eq:one-gadget}
\end{equation}
Both operations preserve the Jacobian determinant and invertibility.  A
collision $p\ne q$ lifts through~\eqref{eq:two-gadget} by appending
$(-P(p),-Q(p))$ and $(-P(q),-Q(q))$, and similarly for
\eqref{eq:one-gadget}.

Starting from~\eqref{eq:phi}, perform the following operations in order.
The symbols $a_j,b_j$ denote the input variables introduced at each step.
\begin{center}
\small
\begin{tabularx}{\textwidth}{@{}c c X c@{}}
\toprule
step & target & new or reused factors & $c$\\
\midrule
1 & $\Phi_1$ & $P=x^2$, $Q=xz+3y$; add $a_1,b_1$ & $-1/2$\\
2 & $\Phi_2$ & $P=3x^2y$, $Q=2z+xyz+3y^2$; add $a_2,b_2$ & $1$\\
3 & $\Phi_2$ & $P=xy$, $Q=a_2z+3xb_2$; add $a_3,b_3$ & $-1$\\
4 & $\Phi_3$ & $P=xy^2$, $Q=7y+3xz+3xy^2+x^2yz$; add $a_4,b_4$ & $1$\\
5 & $\Phi_3$ & add $a_5+a_4xy$; reuse output $b_1+xz+3y$ & $-1$\\
6 & $\Phi_3$ & reuse output $a_3+xy$; add $b_6+a_4b_1-yb_4$ & $1$\\
\bottomrule
\end{tabularx}
\end{center}
Finally postcompose by the shear
\begin{equation}
  M_{b_4}\longmapsto M_{b_4}-M_{a_3}M_{b_1}.
  \label{eq:final-shear}
\end{equation}
The resulting map $M$ has 13 variables
\begin{equation}
 X=(x,y,z,a_1,b_1,a_2,b_2,a_3,b_3,a_4,b_4,a_5,b_6)
 \label{eq:variables}
\end{equation}
and degree three.  Its linear part $L$ is the identity except for
\begin{equation}
 M_{b_1}^{(1)}=b_1+3y,\qquad
 M_{b_2}^{(1)}=b_2+2z,\qquad
 M_{b_4}^{(1)}=b_4+7y.
\end{equation}
In particular $\det L=1$.  Put
\begin{equation}
  \Psi=L^{-1}M=X+H_2(X)+H_3(X),
  \label{eq:psi}
\end{equation}
where $H_d$ is homogeneous of degree $d$.  Then $\deg\Psi=3$,
$\det J\Psi=1$, and $\Psi$ has the following two preimages of the same point:
\begin{align}
p_0={}&(0,0,-\tfrac14,0,0,0,\tfrac12,0,0,0,0,0,0),\nonumber\\
p_1={}&(1,-\tfrac32,\tfrac{13}2,-1,-2,\tfrac92,-10,
\tfrac32,\tfrac34,-\tfrac94,-6,-\tfrac{27}8,\tfrac92),\label{eq:p-collision}\\
\Psi(p_0)=\Psi(p_1)={}&(0,0,-\tfrac14,0,0,0,\tfrac12,0,0,0,0,0,0).
\nonumber
\end{align}
Equations~\eqref{eq:phi} and \eqref{eq:two-gadget}--\eqref{eq:psi} are a
straight-line specification of all components of $\Psi$; no omitted choice is
involved.

\section{The 27-variable cubic homogeneous map}

For $X,Y\in\C^{13}$ and $t\in\C$, define
\begin{equation}
 h(X,Y,t)=\bigl(tH_2(X)+t^2Y,-H_3(X),0\bigr).
 \label{eq:h}
\end{equation}
Every component of $h$ is cubic homogeneous.  Let
$\mathcal B(W)=W+h(W)$.  At $t=1$,
\begin{equation}
 \mathcal B(X,Y,1)=\bigl(X+H_2(X)+Y,Y-H_3(X),1\bigr).
\end{equation}
Consequently, if
\begin{equation}
 r_j=(p_j,H_3(p_j),1),
 \label{eq:rj}
\end{equation}
then $\mathcal B(r_0)=\mathcal B(r_1)=(p_0,0,1)$.  Here $H_3(p_0)=0$, while
\begin{multline}
H_3(p_1)=(-\tfrac{17}4,-\tfrac{45}8,\tfrac{99}{16},0,
\tfrac{135}8,-\tfrac92,-\tfrac{177}8,0,0,\tfrac94,\\
\tfrac{213}8,\tfrac{27}8,0).
\end{multline}

\begin{lemma}\label{lem:nilpotent}
The polynomial matrix $Jh$ is nilpotent.
\end{lemma}
\begin{proof}
Homogeneity and~\eqref{eq:psi} give
\begin{equation}
 \det\bigl(I+tJH_2(X)+t^2JH_3(X)\bigr)
   =\det J\Psi(tX)=1.
 \label{eq:det-homogeneous}
\end{equation}
The Jacobian determinant of $W\mapsto W+h(W)$ is the same determinant by the
block determinant formula.  Hence $\det(I+Jh(W))=1$.  Since $h$ is cubic
homogeneous, $Jh(sW)=s^2Jh(W)$.  Therefore $\det(I+\lambda Jh)=1$
identically in $\lambda$, which is equivalent to nilpotence of $Jh$.
\end{proof}

\section{Symmetrization and the collision}

For a polynomial map $h:\C^r\to\C^r$, de Bondt and van den Essen
\cite{debondt-vandenessen} associate
\begin{equation}
 f_h(A,B)=-\ii\sum_{j=1}^r h_j(A+\ii B)B_j.
\end{equation}
Their characteristic-polynomial identity implies that $\Hess(f_h)$ is
nilpotent if and only if $Jh$ is nilpotent.  Our polynomial in
\eqref{eq:potential} is $\cP=-f_h$, so Lemma~\ref{lem:nilpotent} proves that
$\cP$ is Hessian nilpotent.  It follows immediately that $\det J\Gamma=1$.

For completeness, the collision can be seen directly.  With the linear map
$S(x,y)=(x-\ii y,y)$, differentiation gives
\begin{equation}
 S^{-1}\Gamma S(x,y)=
 \left(x+h(x),(I+Jh(x)^T)y-\ii h(x)\right).
 \label{eq:conjugacy}
\end{equation}
Take $r_0,r_1$ from~\eqref{eq:rj}, let $K_j=I+Jh(r_j)^T$, put
\begin{equation}
 y_1=0,\qquad y_0=K_0^{-1}\bigl(-\ii h(r_1)+\ii h(r_0)\bigr),
\end{equation}
and set $z_j=S(r_j,y_j)$.  Equation~\eqref{eq:conjugacy} and
$r_0+h(r_0)=r_1+h(r_1)$ show $\Gamma(z_0)=\Gamma(z_1)$.  The two 54-tuples
are recorded explicitly in \texttt{collision.json}; all coordinates have
height at most 261.  Direct differentiation of the 598-term expansion checks
this equality exactly over $\Q(\ii)$.

\section{Failure of the Vanishing Conjecture}

Let $\Delta=\sum_j\partial^2/\partial Z_j^2$.  Zhao proved
\cite{zhao} that Hessian nilpotence of a homogeneous polynomial $P$ is
equivalent to
\begin{equation}
  \Delta^mP^m=0\qquad\text{for every }m\geq1.
  \label{eq:zhao-hypothesis}
\end{equation}
Lemma~\ref{lem:nilpotent} and the symmetric reduction therefore show that
$\cP$ satisfies the hypothesis of the Vanishing Conjecture.  Zhao's inversion
formula writes the formal inverse of $Z-t\nabla P$ as $Z+t\nabla Q_t$, where
\begin{equation}
 Q_t=\sum_{m=0}^{\infty}
 \frac{t^m}{2^m m!(m+1)!}\,\Delta^mP^{m+1}.
 \label{eq:zhao-series}
\end{equation}
If $\Delta^m\cP^{m+1}=0$ for all sufficiently large $m$, then $Q_t$ is a
polynomial.  Specializing at $t=1$ would give a polynomial inverse of
$\Gamma$, contradicting the explicit collision above.  Thus
$\Delta^m\cP^{m+1}\ne0$ for infinitely many $m$, proving
Theorem~\ref{thm:main}.

This argument does not exhibit a particular exponent $m$ for which
$\Delta^m\cP^{m+1}\ne0$; it deduces the existence of infinitely many such
exponents indirectly from the collision and the inversion formula.  Zhao's
general finite bound for quartics in $n$ variables reaches
$m>\frac32(3^{n-2}-1)$, or $m>\frac32(3^{52}-1)$ here, so direct expansion at
the general bound is impractical~\cite{zhao}.

\section{Reproducibility}

The repository contains two complementary exact checks.  The script
\texttt{verify.py} rebuilds the construction symbolically and verifies the
original map, all six stable reduction gadgets, the cubic homogeneous map,
Hessian nilpotence through the structural reduction, and the final collision.
The dependency-free script \texttt{verify\_exported\_stdlib.py} uses only
Python's standard library and the two JSON files to verify the expanded
polynomial's degree and term count and the exact collision.  The latter is not
a stand-alone proof of Hessian nilpotence.  All arithmetic in both checks is
exact.  The certificate hashes are
\begin{quote}\small
\texttt{potential\_sparse.json:}\\
\nolinkurl{556683b30fd9e7b9ecfb4c4d5395ee3be6a8dc5f918da08ed60db8c125f0df28}\\[3pt]
\texttt{collision.json:}\\
\nolinkurl{cbea49f88cc59ba7955a4dbd710c68224590ae310a320369c49f8c63e1f78cae}
\end{quote}
The full source and certificates are available in the
\href{https://github.com/AlecKriebel/Math/tree/main/discovery_02_explicit_vanishing_counterexample}{public project repository}.

\section{Scope and novelty}

Zhang's consequence note of 20 July 2026 observes, by Zhao's equivalence, that
the announced three-dimensional counterexample makes the Vanishing Conjecture
false in some finite dimension; it explicitly leaves the consequence
existential~\cite{zhang}.  An earlier public repository~\cite{harrison} gives
a 79-variable cubic homogeneous noninjective Keller map.  Applying the
standard symmetric reduction to that map would produce a 158-variable quartic
in principle, but the repository did not supply an expanded quartic witness
or its transported collision at the time of our search.

The narrower contribution here is an executed 54-variable, 598-term quartic
Hessian-nilpotent polynomial, an exact collision for its gradient map, and a
factor-reusing reduction that keeps the intermediate cubic homogeneous map in
27 variables.  This note does not supply a new proof of the underlying
three-dimensional counterexample or a logically independent disproof of the
Vanishing Conjecture.  To our knowledge, after literature and public-code
searches on 21 July 2026, no earlier source gives this compact explicit
certificate.  We do not claim minimality or unconditional worldwide priority;
the literature surrounding the announced counterexample is evolving rapidly.

\textbf{Post-release correction.} Thompson's initial public
commit~\cite{thompson}, at 03:29:42 UTC on 21 July 2026, already gave a
24-variable cubic homogeneous reduction and therefore beat this note's
27-variable dimension headline before our 13:11:39 UTC release. We failed to
find it in the first audit. Exploration 03 subsequently compressed our
13-variable stable model to a 22-variable cubic and a 44-variable quartic.
Accordingly, this note is superseded for dimension comparisons. Its surviving
role is as the derivation of that stable model and as a reproducible earlier
certificate.

\section*{AI-assistance and verification disclosure}

The construction, proof organization, verification programs, website, and
typeset drafts were developed with extensive assistance from ChatGPT 5.6 Sol.
Alec Kriebel takes responsibility for the submission and for preserving the
complete source and exact certificates.  The algebraic checks are reproducible,
but the note has not been peer reviewed and independent expert scrutiny is
welcome.

\begin{thebibliography}{9}

\bibitem{bcw}
H. Bass, E. H. Connell, and D. Wright,
``The Jacobian conjecture: reduction of degree and formal expansion of the inverse,''
\emph{Bull. Amer. Math. Soc.} 7 (1982), 287--330.
\href{https://doi.org/10.1090/S0273-0979-1982-15032-7}{doi:10.1090/S0273-0979-1982-15032-7}.

\bibitem{debondt-vandenessen}
M. de Bondt and A. van den Essen,
``A reduction of the Jacobian conjecture to the symmetric case,''
\emph{Proc. Amer. Math. Soc.} 133 (2005), 2201--2205.
\href{https://doi.org/10.1090/S0002-9939-05-07570-2}{doi:10.1090/S0002-9939-05-07570-2}.

\bibitem{zhao}
W. Zhao,
``Hessian nilpotent polynomials and the Jacobian conjecture,''
\emph{Trans. Amer. Math. Soc.} 359 (2007), 249--274.
\href{https://arxiv.org/abs/math/0409534}{arXiv:math/0409534}.

\bibitem{alpoge}
L. Alpöge,
X post announcing the three-dimensional map, crediting Akhil Mathew for posing
the question and Claude Fable for producing the example, 20 July 2026.
\href{https://x.com/__alpoge__/status/2079028340955197566}{Original post}.

\bibitem{zhang}
Z. Zhang,
``Direct consequences of the three-dimensional counterexample to the Jacobian
conjecture,'' 20 July 2026.
\href{https://zzhang-iu.github.io/papers/direct-consequences-jacobian/}{Web note}.

\bibitem{harrison}
A. Harrison,
\texttt{jacobian-anatomy}, public GitHub repository, commit
\nolinkurl{74808fb2e1c1691b0007576ba0508e5e7cdcb1e3}, 20 July 2026.
\href{https://github.com/DrAlexHarrison/jacobian-anatomy/commit/74808fb2e1c1691b0007576ba0508e5e7cdcb1e3}{Repository commit}.

\bibitem{thompson}
W. Thompson,
``An explicit 24-variable cubic-homogeneous reduction of the Alp\"oge--Fable
Jacobian counterexample,'' public GitHub repository, commit
\nolinkurl{45a7616fdf5a20c065564f2676190093722696b9}, 21 July 2026.
\href{https://github.com/wtho704/explicit-cubic-homogeneous-jacobian-counterexample}{Repository}.

\end{thebibliography}

\section*{Suggested citation}

\noindent Alec Kriebel, with heavy assistance from ChatGPT 5.6 Sol,
``An explicit Hessian-nilpotent quartic in 54 variables
witnessing the failure of Zhao's Vanishing Conjecture,'' provisional research
note, first posted 21 July 2026.
\url{https://aleckriebel.github.io/Math/papers/explicit-vanishing-counterexample/}.

\end{document}
